\documentclass{article}
\usepackage{amssymb}
\title{The EcoLab Model}
\author{Russell K. Standish}

\newcommand{\br}{\mbox{${\bf r}$}} %reproductive rate
\newcommand{\bbeta}{\mbox{\boldmath{$\beta$}}}   %interaction matrix
\newcommand{\bgamma}{\mbox{\boldmath{$\gamma$}}} %dispersal
\newcommand{\bmu}{\mbox{\boldmath{$\mu$}}}       %mutation rate
\newcommand{\bn}{\mbox{\boldmath{$n$}}}          %species density
\newcommand{\nsp}{\mbox{$n_{\rm sp}$}}           %no. of species           

\providecommand{\reals}{\mathbb{R}}
\providecommand{\naturals}{\mathbb{N}}

\newcommand{\EcoLab}{{\sffamily\slshape
    \mbox{\raisebox{.5ex}{Eco}\hspace{-.4em}{\makebox[.5em]{L}ab}}}}


\begin{document}
\maketitle

\begin{abstract}
  This article documents the EcoLab model, and serves as a repository
  of key behavioural results that may or may not have been published
  in the scientific literature. Earlier versions of this document can
  be found as a section within the official manual of the \EcoLab{}
  platform. It is sometimes referred to as the ``EcoLab Technical Report''.
\end{abstract}

\section{Introduction}

Earlier versions of this paper can be found as a section in the user
manual for the \EcoLab{} simulation system, which provides a reference
implementation of this model.

We start with a generalised form of the Lotka-Volterra equation 
\begin{equation}\label{lotka-volterra}
\dot{\bn} = \br*\bn + \bn*\bbeta\bn.
\end{equation}
Here $\bn\in\reals^n_+$ is the population density, the component
$n_i\ge0$ being the number of individuals of species $i$, \br\ is the
difference between reproduction and death, \bbeta\ is the interaction
matrix, with $\beta_{ij}$ being the interaction between species $i$
and $j$, * referring to elementwise multiplication.

\subsection{Lotka-Volterra Dynamics}

The most obvious thing about equation (\ref{lotka-volterra}) is its
fixed point 
\begin{equation}\label{fixed point}
\hat{\bn} = -\bbeta^{-1}\br,
\end{equation}
where $\dot{\bn}=0$. For this point to be biologically meaningful, all
components of $\hat{\bn}$ must be positive, giving rise to the following
inequalities:
\begin{equation}\label{positive species}
\hat n_i = \left(\bbeta^{-1}\br\right)_i>0, \forall i
\end{equation}
The stability of this point is related to the
negative definiteness of derivative of $\dot{\bn}$ at $\hat{\bn}$. The
components of the derivative are given by
\begin{equation}\label{derivative}
\frac{\partial\dot{n}_i}{\partial n_j} =
\delta_{ij}\left(r_i+\sum_k\beta_{ik}n_k\right) + \beta_{ij}n_i
\end{equation}
Substituting eq (\ref{fixed point}) gives
\begin{equation}
\left.\frac{\partial\dot{n}_i}{\partial n_j}\right|_{\hat{\bn}}=
-\beta_{ij}\left(\bbeta^{-1}\br\right)_i
\end{equation}

Stability of the fixed point requires that this matrix should be
negative definite. Since the $\left(\bbeta^{-1}\br\right)_i$ are
all negative by virtue of (\ref{positive species}), each minor
determinant of this matrix is equal to a minor determinant of \bbeta\
multiplied by a positive number, stability of the equilibrium is
equivalent to \bbeta\ being negative definite.

A weaker condition is to require that the system remain bounded with
time:
\begin{equation}\label{boundedness}
\sum_i\dot{n_i}=\br\cdot\bn + \bn\cdot\bbeta\bn < 0, \;\forall \bn:
\sum_in_i>N \;\exists N
\end{equation}

As \bn\ becomes large in any direction, this functional is dominated
by the quadratic term, so this implies that  $\bn\cdot\bbeta\bn\leq0
\; \forall\bn\in\reals^n_+$. Negative definiteness of \bbeta\ is sufficient,
but not necessary for this condition. Positive eigenvalues of \bbeta\
are possible, provided the associated eigenvectors are not in the
positive cone $\reals^n_+$. For example, the predator-prey
relations (heavily normalised) have the following matrix as \bbeta:
\begin{math}
\bbeta=\left(\begin{array}{cc}
-1 & 2\\
-2 & 0\\
\end{array}\right)
\end{math}
which has eigenvalues $3/2, -5/2$. If we let $\bn=(x,y), x,y\geq0$, then
$\bn\cdot\bbeta\bn=-2x^2$, which is clearly non-positive
for all $x$.

Consider adding a new row and column to \bbeta. What is condition is
the new row and column required to satisfy such that equation
(\ref{boundedness}) is satisfied. Break up \bbeta\ in the following
way:
\begin{displaymath}
\left(
  \mbox{
     \begin{tabular}{c|c}
       $\begin{array}{ccc}\ddots\\&{\bf A}\\&&\ddots\end{array}$ & 
       $\begin{array}{c}\vdots\\{\bf B}\\\vdots\end{array}$ \\
       \hline
       $\begin{array}{ccc}\cdots&{\bf C}&\cdots\end{array}$ & D
     \end{tabular}
   }
\right)
\left(
  \mbox{
     \begin{tabular}{c}
     $\begin{array}{c}\vdots\\{\bf n_1}\\\vdots\end{array}$\\
     \hline
     $n_2$
     \end{tabular}
    }
\right)
\end{displaymath}

Condition (\ref{boundedness}) becomes:
\begin{equation}
{\bf n_1}\cdot{\bf A}{\bf n_1} + {\bf n_1}\cdot({\bf B}+{\bf C})n_2 +
Dn_2^2 \leq 0
\end{equation}

Proceeding inductively, by construction,
${\bf n_1}\cdot{\bf A}{\bf n_1}<0$ and dominate quadratically as
$|{\bf n_1}|\rightarrow\infty$, so it is sufficient to assert, for
$|n_1|=1$:
\begin{equation}\label{boundedness2}
  {\bf n_1}\cdot({\bf B}+{\bf C}) +  Dn_2 \leq 0
\end{equation}

\section{Extinction dynamics}

The original Lotke-Volterra equations (\ref{lotka-volterra}) never
allow for the possibility of extinction, and species density \bn\ can
approach zero arbitrarily closely. In order to simulate actual
extinctions, the EcoLab model uses natural numbers for the density
vector, ie $\bn\in\naturals^n$. Then how to interpret the update
equation (\ref{lotka-volterra})? The EcoLab model is using Euler
integration, with random conversion from real values to the
integers. So, calculate:
\begin{equation}
  \Delta \bn = \br*\bn + \bn*\bbeta\bn
\end{equation}
then for example a value of 4.3 is probabilistically converted to 4
70\% of the time, and to 5 30\% of the time, before being added to the
density vector \bn.

The main alternative scheme for implementing extinction dynamics is to
specify a threshold, under which a species with that density value is
considered extinct. It is unclear what difference this choice of
extinction dynamics makes to the model behaviour, but the
probabilistic approach taken by EcoLab is probably closer to the messy
probabilistic processes found in Nature.

\section{Mutation}\label{mutation}

With mutation, equation (\ref{lotka-volterra}) reads
\begin{equation}
\dot{\bn} = \br*\bn + \bn*\bbeta\bn + {\tt mutate}(\bmu,\br,\bn).
\end{equation}


The difficulty with adding mutation to this model is how to define the
mapping between genotype space and phenotype space, or in other words,
what defines the {\em embryology}. A few studies, including Ray's
Tierra world\cite{Ray91}, do this with an explicit mapping from the
genotype to to some particular organism property (e.g. interpreted as
machine language instructions, or as weight in a neural net). These
organisms then interact with one another to determine the population
dynamics. In this model, however, we are doing away with the
organismal layer, and so an explicit embryology is impossible. The
only possibility left is to use a statistical model of embryology. The
mapping between genotype space and the population parameters $\br$,
$\bbeta{}$ is expected to look like a rugged landscape, however, if
two genotypes are close together (in a Hamming sense) then one might
expect that the phenotypes are likely to be similar, as would the
population parameters. This I call {\em random embryology with
  locality}.

In the simple case of point mutations, the probability $P(x)$ of any
child lying distance $x$ in genotype space from its parent follows a
Poisson distribution. Random embryology with locality implies that the
phenotypic parameters are distributed randomly about the parent
species, with a standard deviation that depends monotonically on the
genotypic displacement. The simplest such model is to distribute the
phenotypic parameters in a Gaussian fashion about the parent's values,
with standard deviation proportional to the genotypic displacement.
This constant of proportionality can be conflated with the species'
intrinsic mutation rate, to give rise another phenotypic parameter
$\bmu$.  It is assumed that the probability of a mutation generating a
previously existing species is negligible, and can be ignored. We also
need another arbitrary parameter $\rho$, ``species radius'', or
\verb+ecolab.sp_sep+,\index{sp\_sep} which can be understood as the
minimum genotypic distance separating species, conflated with the same
constant of proportionality as $\bmu$.

In summary, the mutation algorithm is as follows:
\begin{enumerate}
\item The number of mutant species arising from species $i$ within a
timestep is $\mu_i\alpha_in_i/\rho$. This number is rounded
stochastically to the nearest integer, e.g. 0.25 is rounded up to 1
25\% of the time and down to 0 75\% of the time.

\item Roll a random number from a Poisson distribution
$e^{-x/\mu+\rho}$ to determine the standard deviation $\sigma$ of phenotypic
variation. 

\item Vary $\br$ according to a Gaussian distribution about the
  parents' values, with $\sigma\alpha_0$ as the standard deviation,
  where $\alpha_0$ is the range of values that $\br$ is initialised to,
  ie $\alpha_0$=\verb|ecolab.repro_max|\index{ecolab.repro\_max}$-$
  \verb|ecolab.repro_min|\index{ecolab.repro\_min}

\item The diagonal part of $\bbeta$ must be negative, so vary $\bbeta$
according to a log-normal distribution. This means that if the old
value is $\beta$, the new value becomes
$\beta'=-\exp(\log_e(\beta)+\sigma)$. These values cannot become
arbitrarily small, however, as this would imply that some species make
arbitrarily small demands on the environment, and will become infinite
in number. In EcoLab, the diagonal interactions terms are prevented from
becoming larger than $-r/(.1*{\tt INT\_MAX})$.

\item The offdiagonal components of $\bbeta$, are varied in a similar
fashion to $\br$. However new connections are added, or old ones
removed according to $\lfloor 1/r\rfloor$, where $r\in(-2,2)$ is
chosen from a stepped uniform distribution 
\begin{displaymath}
P(x)=\left\{
\begin{array}{ll}
0.25(1-g) & \mathrm{if}\; x\leq0\\
0.25(1+g) & \mathrm{if}\; x>0\\
\end{array}
\right.
\end{displaymath}
where $g\in[-1,1]$ (default of 0) is specified by 
\verb+ecolab.gen_bias+\index{generalization\_bias}. The values on
the new connections are chosen from the same initial distribution that
the offdiagonal values where originally set with, ie the range
\verb|ecolab.odiag_min|\index{ecolab.odiag\_min} to
\verb|ecolab.odiag_max|\index{ecolab.odiag\_max}. 

To convert  condition (\ref{boundedness2}) into something tractable for the mutation
operator, we make use of the fact that $n\in\naturals^n$.

Let 
\begin{displaymath}
b=\sum_{i: B_i+C_i>0}B_i+C_i.
\end{displaymath}
Then a sufficient but
not necessary condition for condition (\ref{boundedness2}) is
\begin{displaymath}
b+Dn_2\leq0
\end{displaymath}

Since $D<0$, the maximum of $b+Dn_2$ for
$n_2\in\naturals$ occurs for $n_2=1$, so the condition becomes:
\begin{equation}\label{final-condition}
  D<-b.
\end{equation}

Originally, this condition was achieved by adjusting the values of the
offdiagonal components added to the interaction matrix \bbeta, but
this had the unfortunate effect of gradually reducing the connectivity
over simulation time, leading to a natural growth in diversity (see
\S\ref{nsp-connectivity}). So now, we adjust the newly added diagonal
term $D$ to be strictly less than $\max\{-b, 0\}$, retaining a natural
conservation of connectivity under mutation dynamics.

\item $\bmu$ must be positive, so should evolve according to the
  log-normal distribution like the diagonal components of $\bbeta$.
  Similar to $\bbeta$, it is a catastrophe to allow $\bmu$ to become
  arbitrarily large. In the real world, mutation normally exists at
  some fixed background rate --- species can reduce the level of
  mutation by improving their genetic repair algorithms. In the EcoLab
  model, this ceiling on $\bmu$ is given by the
  \verb|ecolab.mut_max|\index{ecolab.mut\_max} variable.

\end{enumerate}

\section{Criticality}

Robert May\cite{May72} noted that the phase space spanned by
connectivity
\begin{equation}
\sigma^2=\sum_{i,j: i\ne j}
\frac{\beta_{ij}^2}{\beta_{ii}\beta_{jj}n_\mathrm{sp}(n_\mathrm{sp}-1)}
\end{equation}
and diversity
$n_\mathrm{sp}$ had a hyperbolic critical line
\begin{equation}
  n_\mathrm{sp}=1/\sigma^2
\end{equation}
under which the fixed point $\hat{\bn}$ is
most likely to be stable, and above which it is most likely to be
unstable.

Stability of the fixed is not really the right concept with these
dynamics, after all the classic Lotka-Volterra systems has an unstable
fixed point, but a limit cycle instead. Systems with more than 3
dimension can have strange attractors as well.

Persistence is a more relevant concept. Persistence is the property
that no species in the ecosystem goes extinct under conditions of no
mutation ($\bmu=0$). 

However, the EcoLab model does exhibit a similar critical

\section{Input Parameters}\label{input parameters}

The model's parameters are set in the python scripts.
The actual data structures of the model are initialised the first time
the model's generate step is called. An example input set is:
\begin{verbatim}
from ecolab_model import panmictic_ecolab as ecolab
ecolab.species([1,2])
ecolab.density([100,100]) 
ecolab.create([0,0])
ecolab.repro_rate([.1,-.1])
ecolab.interaction.diag([-.0001,-1e-5])
ecolab.interaction.val([-0.001,0.001])
ecolab.interaction.row([0,1])
ecolab.interaction.col([1,0])
ecolab.mutation([.01,.01])
ecolab.sp_sep(.1)
ecolab.repro_min(-.1)
ecolab.repro_max(.1)
ecolab.odiag_min(-1e-3)
ecolab.odiag_max(1e-3)
ecolab.mut_max(.01)
\end{verbatim}

Model variables define a python method of the same name as they appear
in the C++ source. So in the EcoLab model, the C++ object {\tt
  ecolab} defines a set of python methods such as {\tt ecolab.density}
that can be used for setting or querying the values of {\tt ecolab}'s
members.  If an argument is specified, then that argument is used to
set the variable's value, otherwise, the variable's value is returned.
Array members in the model are initialised by specifying a python
array argument to the variables name, and return python objects when
no argument is specified that can be converted to python arrays by
calling them. For example \verb+ecolab.density()()+ returns the
current density array. The above example starts the ecology off with a
single predator and prey (based on {\tt
  pred-prey.py}\index{pred-prey.py}).

\section{EcoLab Model commands}

Archetypal EcoLab experimental scripts can be found at
\verb+panmictic_ecolab.py+ for no spatial variance, and
\verb+spatial_ecolab.py+ for the spatial model. EcoLab provides a
static object for each type: \verb+panmictic_ecolab+ and \verb+spatial_ecolab+.

\subsection{generate}\index{generate}

This implements the basic Lotka-Volterra equations:

\begin{math}
\dot{\bn} = \br*\bn + \bn*\bbeta\bn
\end{math}

with \br\ being the reproduction rate and \bbeta\ being the
interspecies interaction. This can be implemented as a single line:
\begin{verbatim}
  density += repro_rate * density + (interaction * density) * density;
\end{verbatim}
although for performance reasons on GPUs, the implementation is more
complicated.

This command also increments the timestep counter {\tt tstep}.

An optional argument specifies a number of timesteps to run the
generate step. This improves the speed by amortising the real to
integer conversion operation over a number of timesteps. The downside
is that computation may fail if the problem is ill-conditioned
(offdiagonal elements of \bbeta\ too large with respect to the
diagonal elements).

\subsection{condense}\label{condense}\index{condense}

This command compacts the systems of equations by removing extinct
species where $n_i=0$. It returns the number of extinctions.

\subsection{mutate}\label{mutate}\index{mutate}

This applies the point mutations to the system. The precise algorithm
is described in \S\ref{mutation}. 

\subsection{migrate}\label{migrate}\index{migrate}

This operator implements migration within cellular \EcoLab{}. This
updates density values according to the difference with the 4 nearest
neighbours: $\bn=\bn+\bgamma*(0.25(\bn_n+\bn_e+\bn_s+\bn_w)-\bn)$, where the
$n,e,s,w$ index the north, east, south and west neigbouring cells.

\subsection{lifetimes}\label{lifetimes}\index{lifetimes}

\verb|lifetimes| records the timestep when a species is created in the
{\tt create}\index{create} iarray. If a species has gone extinct, the
value in {\tt create} is zero. Upon return, this routine returns the
lifetime of the species that have gone extinct. This can then be
passed to a histogram routine, or written to a file.

\subsection{random\_interaction(conn, sigma)}\label{random_interaction}
\index{random\_interaction}

Calls {\protect\tt sparse\_mat::init\_rand()} (\S\ref{sparse_mat}) to randomly initialising the nonzero pattern of
the offdiagonal elements. The average number of nonzeros per row is
{\tt conn}, and the standard deviation of the number of nonzeros is
{\tt sigma}.

\subsection{seed({\em seed})}

Sets the global random number seed used for initialisation and
rounding floating values to integers, eg in the {\bf generate} step.

\subsection{complexity}

Returns the complexity of the foodweb, using the algorithm described
in \cite{Standish10b}.

\subsection{connectivity $\sigma^2$}

Returns the connectivity of the interaction matrix, defined
analogously to \cite{May72} as the variance of offdiagonal components
of the interaction matrix \bbeta, normalised so that the diagonal
coefficients are unity.

\begin{equation}
  \sigma^2=\sum_{i,j: i\ne j} \frac{\beta_{ij}^2}{\beta_{ii}\beta_{jj}n_\mathrm{sp}(n_\mathrm{sp}-1)}
\end{equation}
\subsection{setGrid}

Synopsis

{\tt ecolab.setGrid}({\em x}, {\em y})

\noindent Set up an $x\times y$ grid in the spatial EcoLab model. See
\verb+spatial_ecolab.py+ for an example using this.

\subsubsection{cell}

Synopsis

{\tt ecolab.cell}({\em x}, {\em y})

\noindent Creates a python reference to the cell located at $x$ and
$y$. For example
\begin{verbatim}
c=ecolab.cell(1,0)
print(c.density()())
\end{verbatim}

\subsection{nsp}

Returns the array of diversities (number of species) of each cell.

\section{ConnectionPlot}

The connectionPlot member of the panmictic model visualises the
interaction matrix, with each element coloured by the connected group
it belongs to. Extinct species are coloured in buff. To display the
contents requires a little setup, but is an example of how to
implement custom visualisers in \EcoLab.

\begin{verbatim}
connectionPlot=Tk()
connectionPlot.wm_title('Connections')
connectionPlot.eval('load '+ecolabHome()+'/lib/ecolab.so')
connectionPlot.eval('image create cairoSurface connectionCanvas '+
   '-surface ecolab_model panmictic_ecolab.connectionPlot')
ttk.Label(connectionPlot, image='connectionCanvas').
  pack(fill='both',expand=True)
windows.append(ref(connectionPlot))
ecolab.updateConnectionPlot()
\end{verbatim}

Subsequent calls to \verb+updateConnectionPlot+ redraws the image from
the current state of the interaction marix \bbeta.

\section{Checkpoint/restart}

The EcoLab model, like all \EcoLab{} models, supports long running
simulations with the ability to write a checkpoint, then later restart
the simulation from that checkpoint. The \verb+checkpoint+ and
\verb+restart+ commands take a filename to write or read the
checkpoint.

\section{Spatial Variation}\label{spatial}

The ecolab model can run in multicellular mode by importing the
\verb+spatial_ecolab+ static variable into Python, and calling
\verb+setGrid+ to set the dimensions of the grid:
\begin{verbatim}
from ecolab_model import spatial_ecolab as ecolab
numX=8
numY=8
ecolab.setGrid(numX,numY)
\end{verbatim}

Only population density varies between the cells --- all other
variables are members of the ecolab variable so can be set or queried
in the usual way.

The usual ecolab model methods (generate, mutate, condense and
lifetimes)\index{generate}\index{mutate}\index{condense}\index{lifetimes}
can now be called, but operate on the entire grid. A new {\tt
migrate}\index{migrate} method is defined to handle migration between
cells.

Access to the individual cells can obtained by calling \verb+cell+ on
the SpatialEcolab object.

\subsection{Parallel Execution}

Since the cells are pins in a graphcode Graph (\S\ref{}{graphcode}),
they are distributed over parallel processes if available.

Remember to call the \verb+gather+\index{gather} method to ensure node
0 is updated before running a visualiser on global data.

\begin{description}
\item[ecolab.gather] Bring processor 0's data up to date with the rest
  of the grid
\item[ecolab.distributeObjects] Broadcast processor 0's data out to the
  rest of the grid.
\end{description}

\end{document}
